
\appendix
\section{}
\section{Properties of correlation functions} 
\label{appprop}
%----------------------------------------------------------------------

In this section we prove the theorems stated in section \ref{propcorrel}.


We use very much the following obvious properties:

\beq\label{sumBdxdx}
\sum_i B(p^i,q) = {dx(p)dx(q)\over (x(p)-x(q))^2},
\eeq
\beq\label{eqsymdEom}
{dE_{\qbar}(p)\over \om(\qbar)} = {dE_{q}(p)\over \om(q)},
\eeq
and
\beq\label{eqdEtoB}
\mathop{{\rm lim}}_{q\to a_i} {dE_{q}(p)\over \om(q)}\,dx(q) = - {1\over 2}\,\mathop{{\rm lim}}_{q\to a_i} {B(p,q)\over dy(q)}\,dx(q).
\eeq
The third one is nothing but De L'H\^opital's rule.

\vs

%---------------------------------thpolesW30------------------
{\bf Theorem \ref{thW30} }\proof{
From the definition \ref{defloopfctions}, we have:
\bea
W_{3}^{(0)}(p,p_1,p_2)
&=&  \Res_{q\to \bfa} {dE_{q}(p)\over \om(q)}\,\Big(B(q,p_1)B(\qbar,p_2)+B(q,p_2)B(\qbar,p_1) \Big) \cr
&=& 2 \Res_{q\to \bfa} {dE_{q}(p)\over \om(q)}\,B(q,p_1)B(\qbar,p_2) \cr
&=& 2 \Res_{q\to \bfa} \Res_{r\to \overline{q}} {dE_{q}(p)\over (y(q)-y(\qbar))(x(r)-x(q))}\,B(q,p_1)B(r,p_2) \cr
&=& - 2 \Res_{q\to \bfa} \Res_{r\to q} {dE_{q}(p)\over (y(q)-y(\qbar))(x(r)-x(q))}\,B(q,p_1)B(r,p_2) \cr
&=& - 2 \Res_{q\to \bfa} {dE_{q}(p)\over \om(q)}\,B(q,p_1)B(q,p_2) \cr
&=& \Res_{q\to \bfa} {B(q,p)B(q,p_1)B(q,p_2)\over dx(q) dy(q)}
\eea
where we have used \eq{eqsymdEom}, \eq{sumBdxdx} and \eq{eqdEtoB}.
}

\vs

%---------------------------------thpolesWkgbp------------------
{\bf Theorem \ref{thpolesWkgbp}} \proof{
It is  obvious from the definition that if $p$ is away from branchpoints, the residues are finite integrals, and $W_{k+1}^{(g)}$ is finite.
The only poles can be obtained when $p$ pinches an integration contour, i.e. at branch points.

Then, it is easy to see by recursion on $k$ and $g$ that it holds for any $p_i$.

}
\vs

%---------------------------------thWkcycle------------------
{\bf Theorem \ref{thWkcycle}} \proof{The first one is a property of $dE_{q}(p)$, and the second one follows from recursion on $k$ and $g$, and it holds for $W_{2}^{(0)}=B$.}

\vs

%---------------------------------thsumWk------------------
{\bf Theorem \ref{thsumWk}} \proof{
The case $k=1,g=0$ comes from \eq{sumBdxdx}, and by integration:
\beq
\sum_i dE_{q}(p^i) = 0
\eeq
which proves equation \ref{thsumWk1}.
Then, it is clear from the iterative definition of $W_{k+1}^{(g)}$ for $k\geq 1$,
that the dependance of $W_{k+1}^{(g)}(p_1,p,p_2,\dots,p_k)$  in $p$ is a sum of integrals involving only the following quantities:
\beq
\sum_j \Res_{q\to a_j}\,{dE_{j,q}(q_1)\over \om_j(q)} B(q,p)\, W_{l+1}^{(g)}(q, q_2,\dots,q_l)
\eeq
for some $q_1,\dots,q_l$\footnote{Since it may not be clear which branch point $dE$ and $\omega$ refers to,
one indicates it with an index.}.
Using equation \ref{sumBdxdx} we have to compute:
\bea
&& \sum_i \sum_j \Res_{q\to a_j}\,{dE_{j,q}(q_1)\over \om_j(q)} B(q,p^i)\, W_{l+1}^{(g)}(q, q_2,\dots,q_l) \cr
&=&  \sum_j \Res_{q\to a_j}\,{dE_{j,q}(q_1)\over \om_j(q)} {dx(q)dx(p)\over (x(q)-x(p))^2}\, W_{l+1}^{(g)}(q, q_2,\dots,q_l) \cr
&=& {1\over 2}\, \sum_j \Res_{q\to a_j}\,{dE_{j,q}(q_1)\over \om_j(q)} {dx(q)dx(p)\over (x(q)-x(p))^2}\, \left( W_{l+1}^{(g)}(q, q_2,\dots,q_l)+W_{l+1}^{(g)}(\qbar, q_2,\dots,q_l)\right) \cr
&=& - {1\over 2}\,\sum_{q^i\neq q,\qbar}\, \sum_j \Res_{q\to a_j}\,{dE_{j,q}(q_1)\over \om_j(q)} {dx(q)dx(p)\over (x(q)-x(p))^2}\,  W_{l+1}^{(g)}(q^i, q_2,\dots,q_l) \cr
&=& 0
\eea
where the second equality holds due to \ref{eqsymdEom}, the third equality holds due to equation \ref{thsumWk1}, and the last equality holds because that last expression has no poles at the branchpoints.
This proves equation \ref{thsumWk2}.
}

\vs

%---------------------------------thPkpol------------------
{\bf Theorem \ref{thPkpol}} \proof{
It is clearly a rational function of $x(p)$ because it is a symmetric sum on all sheets.
From theorem \ref{thpolesWkgbp}, the RHS may have poles at branch points, and/or at the poles of some $y(p^i)$, and/or when $x(p)=x(p_l)$ for some $l$.
Let us prove that the poles at branch points actualy cancel.

Let us denote in this section:
\beq
U_k^{(g)}(q,q',p_K) =
\sum_{m=0}^g \sum_{J\subset K} W_{j+1}^{(m)}(q,p_J)W_{k-j+1}^{(g-m)}(q',p_{K/J})
+ W_{k+2}^{(g-1)}(q,q',p_K).
\eeq
From theorem \ref{thsumWk}, we have:
\bea
&& \sum_i U_k^{(g)}(q^i,q^i,p_K) \cr
&=& -\sum_i\sum_{l\neq i} U_k^{(g)}(q^l,q^i,p_K) \cr
&=& - U_k^{(g)}(q,\qbar,p_K) - U_k^{(g)}(\qbar,q,p_K) - \sum_{q^i\neq q,\qbar} (U_k^{(g)}(q,q^i,p_K)+U_k^{(g)}(\qbar,q^i,p_K)) \cr
&& - \sum_{q^i\neq q,\qbar} (U_k^{(g)}(q^i,q,p_K)+U_k^{(g)}(q^i,\qbar,p_K))  - \sum_{q^i\neq q,\qbar\,}\,\, \sum_{l\neq i, q^l\neq, q,\qbar} U_k^{(g)}(q^l,q^i,p_K) \cr
\eea
and from theorem \ref{thsumWk} and theorem \ref{thpolesWkgbp}, only the first two terms have poles at branch points, and thus:
\bea
&& - \sum_i \sum_j \Res_{q\to a_j} {dE_{j,q}(p)\over \om_j(q)}\, U_k^{(g)}(q^i,q^i,p_K) \cr
&=&  \sum_j \Res_{q\to a_j} {dE_{j,q}(p)\over \om_j(q)}\, \left(U_k^{(g)}(q,\qbar,p_K)+U_k^{(g)}(\qbar,q,p_K)\right) \cr
&=&  2 \sum_j \Res_{q\to a_j} {dE_{j,q}(p)\over \om_j(q)}\, U_k^{(g)}(q,\qbar,p_K) \cr
&=& 2 W_{k+1}^{(g)}(p,p_K)
\eea
where the last equality holds from the definition of $W_{k+1}^{(g)}$.
Thus we have:
\bea
W_{k+1}^{(g)}(p,p_K)
&=& - {1\over 2}\,\sum_j \Res_{q\to a_j} {dE_{j,q}(p)\over \om_j(q)}\, \left( \sum_i U_k^{(g)}(q^i,q^i,p_K)  \right). \cr
\eea
Then we rewrite it in terms of $P_k^{(g)}$ above:
\bea
&& W_{k+1}^{(g)}(p,p_K) \cr
&=& - {1\over 2}\,\sum_j \Res_{q\to a_j} {dE_{j,q}(p)\over \om_j(q)}\, \Big( P_k^{(g)}(x(q),p_K)\,dx(q)^2  \cr
&& \qquad \quad \qquad + 2\sum_i y(q^i)dx(q)W_{k+1}^{(g)}(q^i,p_K)  \Big) \cr
&=& - {1\over 2}\,\sum_j \Res_{q\to a_j} {dE_{j,q}(p)\over \om_j(q)}\, \Big( P_k^{(g)}(x(q),p_K)\,dx(q)^2  + 2y(q)dx(q)W_{k+1}^{(g)}(q,p_K) \cr
&&  \qquad \quad \qquad + 2y(\qbar)dx(q)W_{k+1}^{(g)}(\qbar,p_K) \Big) \cr
&=& - {1\over 2}\,\sum_j \Res_{q\to a_j} {dE_{j,q}(p)\over \om_j(q)}\, \Big( P_k^{(g)}(x(q),p_K)\,dx(q)^2  \cr
&& \qquad \quad \qquad + 2(y(q)-y(\qbar))dx(q)W_{k+1}^{(g)}(q,p_K) \Big) \cr
\eea
where we have used theorem \ref{thsumWk} again. That gives
\bea
&& W_{k+1}^{(g)}(p,p_K) \cr
&=& - {1\over 2}\,\sum_j \Res_{q\to a_j} {dE_{j,q}(p)\over \om_j(q)}\,  P_k^{(g)}(x(q),p_K)\,dx(q)^2 \cr
&& \qquad \quad - \sum_j \Res_{q\to a_j} dE_{j,q}(p)\,W_{k+1}^{(g)}(q,p_K).  \cr
\eea
Let us compute that last integral:
\bea\label{proofthPkpolRbilin}
&& \sum_j \Res_{q\to a_j} dE_{j,q}(p)\,W_{k+1}^{(g)}(q,p_K)  \cr
&=& - \Res_{q\to p} dE_{j,q}(p)\,W_{k+1}^{(g)}(q,p_K) + {1\over 2i\pi}\,\sum_i \oint_{q'\in\underline\bcycle_i}B(p,q')\, \oint_{\underline\acycle_i} W_{k+1}^{(g)}(q,p_K)  \cr
&& -  {1\over 2i\pi}\,\sum_i \oint_{{q'\in\underline\acycle_i}}B(p,q')\, \oint_{\underline\bcycle_i} W_{k+1}^{(g)}(q,p_K)  \cr
&=& - \Res_{q\to p} dE_{j,q}(p)\,W_{k+1}^{(g)}(q,p_K) + \sum_i  du_i(p)\, \oint_{\acycle_i} W_{k+1}^{(g)}(q,p_K)  \cr
&=& - \Res_{q\to p} dE_{j,q}(p)\,W_{k+1}^{(g)}(q,p_K)   \cr
&=& - W_{k+1}^{(g)}(p,p_K)
\eea
where we have deformed the contour of integration using Riemann bilinear identity, and then we have used theorem \ref{thWkcycle}.
Thus we have:
\beq
\sum_j \Res_{q\to a_j} {dE_{j,q}(p)\over \om_j(q)}\,  P_k^{(g)}(x(q),p_K)\,dx(q)^2  = 0 .
\eeq
Since this holds for any $p$, we can write for any $m\geq 0$:
\bea
0
&=& \Res_{p\to a_i} (y(p)-y(a_i))\,(x(p)-x(a_i))^m\, \sum_j \Res_{q\to a_j} {dE_{j,q}(p)\over \om_j(q)}\,  P_k^{(g)}(x(q),p_K)\,dx(q)^2  \cr
&=& - \Res_{q\to a_i}   \Res_{p\to q,\qbar} {dE_{j,q}(p)\over \om_j(q)}\,  (y(p)-y(a_i))\,(x(p)-x(a_i))^m\,P_k^{(g)}(x(q),p_K)\,dx(q)^2  \cr
&=& - \Res_{q\to a_i} dx(q)\,  (x(q)-x(a_i))^m\,P_k^{(g)}(x(q),p_K)  \cr
\eea
which proves that $P_k^{(g)}(x(q),p_K)$ can have no pole at $q=a_i$.

}

\vs

%---------------------------------thsymWk------------------
{\bf Theorem \ref{thsymWk}} \proof{
Assume this is proved for $h<g$, and at $g$, it is proved for $l<k$.
Then we have:
\bea
&& W_{2}^{(g)}(p_1,p_2) \cr
&=&  \sum_i \Res_{q\to a_i} {dE_{q}(p_1)\over \om(q)}\,\Big[2\sum_{m=0}^g  W_{2}^{(m)}(q,p_2)W_{1}^{(g-m)}(\qbar)
+ W_{3}^{(g-1)}(q,\qbar,p_2) \Big] \cr
&=&  \sum_i \Res_{q\to a_i} {dE_{q}(p_1)\over \om(q)}\,B(q,p_2)\,W_{1}^{(g)}(\qbar) \cr
&& + \sum_i \Res_{q\to a_i} \sum_{i'} \Res_{q'\to a_{i'}} {dE_{q}(p_1)\over \om(q)}\,{dE_{i',q'}(p_2)\over \om_{i'}(q')}\,  \Big[ 2  \sum_{m'=0}^{g-1}  W_{2}^{(m')}(q',q)W_{2}^{(g-m'-1)}(\qbar',\qbar) \cr
&& \quad + 2 \sum_{m=1}^g W_{1}^{(g-m)}(\qbar)W_{3}^{(m-1)}(q',\qbar',q)  +  2  \sum_{m'=0}^{g-1}  W_{3}^{(m')}(q',q,\qbar)W_{1}^{(g-m'-1)}(\qbar') \cr
&& \qquad + 4\sum_{m=1}^g\sum_{m'=0}^m  W_{2}^{(m')}(q',q)W_{1}^{(m-m')}(\qbar')W_{1}^{(g-m)}(\qbar) + W_{4}^{(g-2)}(q,q',\qbar,\qbar') \Big] \cr
\eea
and $W_{2}^{(g)}(p_2,p_1)$ is given by the same integral except that the order for computing residues is reversed, the residue in $q$ is computed before $q'$.
The difference is thus obtained by pushing the contour of $q'$ through the contour of $q$, and is obtained as the residue at $q=q'$.
Notice that only $W_2^{(0)}(q,q')$ has a pole at $q=q'$, all the other $W_k^{(g)}$ have no poles at $q=q'$ from theorem \ref{thpolesWkgbp}. Thus we have:
\beq
\begin{array}{l}
W_{2}^{(g)}(p_1,p_2)-W_{2}^{(g)}(p_2,p_1)
- \sum_i \Res_{q\to a_i} {dE_{q}(p_1)\,B(q,p_2)-dE_{q}(p_2)\,B(q,p_1)\over \om(q)}\,W_{1}^{(g)}(\qbar) \cr
=  \sum_i \Res_{q\to a_i} \Res_{q'\to q} {dE_{q}(p_1)\over \om(q)}\,{dE_{q'}(p_2)\over \om_{i}(q')}\, B(q,q') \cr
 \qquad \Big[4\sum_{m=0}^g  W_{1}^{(m)}(\qbar')W_{1}^{(g-m)}(\qbar)  + 2  W_{2}^{(g-1)}(\qbar',\qbar)  \Big] \cr
= 2 \sum_i \Res_{q\to a_i} \Res_{q'\to q} {dE_{q}(p_1)\over \om(q)}\,{dE_{q'}(p_2)\over \om_{i}(q')}\, B(q,q')\, U_{0}^{(g)}(\qbar,\qbar') \cr
=  \sum_i \Res_{q\to a_i} \Res_{q'\to q} {dE_{q}(p_1)\over \om(q)}\,{dE_{q'}(p_2)\over \om_{i}(q')}\, B(q,q')\, \left( U_{0}^{(g)}(q,q') + U_{0}^{(g)}(\qbar,\qbar')\right) \cr
\end{array}\eeq
where we have used the notation of theorem \ref{thPkpol}.

Similarly, for higher values of $k$, we find:
\bea
&& W_{2+k}^{(g)}(p_1,p_2,p_K)-W_{2+k}^{(g)}(p_2,p_1,p_K) \cr
&& \qquad - \sum_i \Res_{q\to a_i} {dE_{q}(p_1)\,B(q,p_2)-dE_{q}(p_2)\,B(q,p_1)\over \om(q)}\,W_{k+1}^{(g)}(\qbar,p_K) \cr
&=&  \sum_i \Res_{q\to a_i} \Res_{q'\to q} {dE_{q}(p_1)\over \om(q)}\,{dE_{q'}(p_2)\over \om_{i}(q')}\, B(q,q') \,\left(U_{k}^{(g)}(q,q',p_K)+U_{k}^{(g)}(\qbar,\qbar',p_K)\right) \cr
\eea
Then it gives:
\bea
&& W_{2+k}^{(g)}(p_1,p_2,p_K)-W_{2+k}^{(g)}(p_2,p_1,p_K) \cr
&& \qquad - \sum_i \Res_{q\to a_i} {dE_{q}(p_1)\,B(q,p_2)-dE_{q}(p_2)\,B(q,p_1)\over \om(q)}\,W_{k+1}^{(g)}(\qbar,p_K) \cr
&=&  \sum_i \Res_{q\to a_i} {dE_{q}(p_1)\over \om(q)}\,d_{q'}\left({dE_{q'}(p_2)\over \om_{i}(q')}\, \left(U_{k}^{(g)}(q,q',p_K)+U_{k}^{(g)}(\qbar,\qbar',p_K)\right)\right)_{q'=q} \cr
\eea
and by integrating half of it by parts we get:
\beq
\begin{array}{l}
 W_{2+k}^{(g)}(p_1,p_2,p_K)-W_{2+k}^{(g)}(p_2,p_1,p_K) \cr
= - \sum_i \Res_{q\to a_i} {B(q,p_1) dE_{q}(p_2)-B(q,p_2) dE_{q}(p_1)\over \om(q)^2}\, \left(U_{k}^{(g)}(q,q,p_K)+U_{k}^{(g)}(\qbar,\qbar,p_K)\right) \cr
 +  \sum_i \Res_{q\to a_i} {dE_{q}(p_1)\,B(q,p_2)-dE_{q}(p_2)\,B(q,p_1)\over \om(q)}\,W_{k+1}^{(g)}(\qbar,p_K) \cr
\end{array}
\eeq

Now we use the following Lemma:
\bl\label{LemmathsymWk}
If $f(q,q')$ is localy a bilinear differential near a branchpoint $a_i$, with no poles, and symmetric in $q$ and $\qbar$, and in $q'$ and $\qbar'$, then:
\beq
\Res_{q\to a_i} {B(q,p_1) dE_{q}(p_2)-B(q,p_2) dE_{q}(p_1)\over \om(q)^2} f(q,q) =0.
\eeq
\el
{\bf Proof of the Lemma:}

The residue is a simple pole, and we can use formula \ref{eqdEtoB}, that gives:
\bea
&& \Res_{q\to a_i} {B(q,p_1) dE_{q}(p_2)-B(q,p_2) dE_{q}(p_1)\over \om(q)^2} f(q,q) \cr
&=& \Res_{q\to a_i} {B(q,p_1) B(q,p_2)-B(q,p_2) B(q,p_1)\over \om(q) dy(q) dx(q)} f(q,q) \cr
&=& 0
\eea
$\square$.
\medskip

Using this Lemma, as well as theorem \ref{thPkpol}, we get:
\beq
\begin{array}{l}
W_{2+k}^{(g)}(p_1,p_2,p_K)-W_{2+k}^{(g)}(p_2,p_1,p_K) \cr
= - \sum_i \Res_{q\to a_i} {B(q,p_1) dE_{q}(p_2)-B(q,p_2) dE_{q}(p_1)\over \om(q)^2}\, (y(q)-y(\qbar))dx(q)W_{k+1}^{(g)}(q,p_K) \cr
 +  \sum_i \Res_{q\to a_i} {dE_{q}(p_1)\,B(q,p_2)-dE_{q}(p_2)\,B(q,p_1)\over \om(q)}\,W_{k+1}^{(g)}(\qbar,p_K) \cr
= 0.\cr
\end{array}
\eeq

}

\vs

{\bf Corollary \ref{corResxyWk}} \proof{

For any rational function $R(x)$ with no pole at $x(a_i)$ we have:
\bea
&& \Res_{a_i} R(x(p)) W^{(g)}_{k+1}(p,p_1,\dots,p_k)  \cr
&=& {1\over 2} \Res_{a_i} R(x(p)) \left(W^{(g)}_{k+1}(p,p_1,\dots,p_k)+W^{(g)}_{k+1}(\pbar,p_1,\dots,p_k)\right)  \cr
&=& 0
\eea
due to theorem \ref{thsumWk}.

\medskip

For $m=0,1$ compute
\bea
&& \sum_\alpha \Res_{p\to \alpha} x(p)^m y(p) W^{(g)}_{k+1}(p,p_K) \cr
&=& -{1\over 2}\sum_\alpha \Res_{p\to \alpha} {x(p)^m\over dx(p)}\, \Big( -2y(p)dx(p) W^{(g)}_{k+1}(p,p_K)+ \cr
&& \qquad +\sum_{h=0}^g \sum_{I\subset K} W^{(h)}_{|I|+1}(p,p_I)W^{(g-h)}_{k-|I|+1}(p,p_{K/I})+W^{(g-1)}_{k+2}(p,p,p_K)\Big)\cr
&=& {1\over 2} \Res_{p\to a_i,p_K} {x(p)^m\over dx(p)}\, \Big( -2y(p)dx(p) W^{(g)}_{k+1}(p,p_K)+ \cr
&& \qquad +\sum_{h=0}^g \sum_{I\subset K} W^{(h)}_{|I|+1}(p,p_I)W^{(g-h)}_{k-|I|+1}(p,p_{K/I})+W^{(g-1)}_{k+2}(p,p,p_K)\Big)\cr
&=& {1\over 2}\sum_{j=1}^k \Res_{p\to p_j} {x(p)^m\over dx(p)}\, \Big( B(p,p_j) W^{(g)}_{k}(p,p_{K/\{j\}})\Big)\cr
&& +{1\over 2}\sum_i \Res_{p\to a_i} {x(p)^m\over dx(p)}\, \Big( -2y(p)dx(p) W^{(g)}_{k+1}(p,p_K)+ \cr
&& \qquad +\sum_{h=0}^g \sum_{I\subset K} W^{(h)}_{|I|+1}(p,p_I)W^{(g-h)}_{k-|I|+1}(p,p_{K/I})+W^{(g-1)}_{k+2}(p,p,p_K)\Big)\cr
&=& {1\over 2}\sum_{j=1}^k d_{p_j}\, \Big( {x(p_j)^m\,W^{(g)}_{k}(p_{K})\over dx(p_j)}  \Big)\cr
&& +{1\over 4}\sum_i \Res_{p\to a_i} {x(p)^m\over dx(p)}\, \Big( -2y(p)dx(p) W^{(g)}_{k+1}(p,p_K)+ \cr
&& \qquad +\sum_{h=0}^g \sum_{I\subset K} W^{(h)}_{|I|+1}(p,p_I)W^{(g-h)}_{k-|I|+1}(p,p_{K/I})+W^{(g-1)}_{k+2}(p,p,p_K)\Big)\cr
&& +{1\over 4}\sum_i \Res_{p\to a_i} {x(p)^m\over dx(p)}\, \Big( -2y(\pbar)dx(p) W^{(g)}_{k+1}(\pbar,p_K)+ \cr
&& \qquad +\sum_{h=0}^g \sum_{I\subset K} W^{(h)}_{|I|+1}(\pbar,p_I)W^{(g-h)}_{k-|I|+1}(\pbar,p_{K/I})+W^{(g-1)}_{k+2}(\pbar,\pbar,p_K)\Big)\cr
&=& {1\over 2}\sum_{j=1}^k d_{p_j}\, \Big( {x(p_j)^m\,W^{(g)}_{k}(p_{K})\over dx(p_j)}  \Big)\cr
&& +{1\over 4}\sum_i \Res_{p\to a_i} {x(p)^m\over dx(p)}\, \Big( P_k^{(g)}(x(p),p_K) dx^2(p) \Big) \cr
&=& {1\over 2}\sum_{j=1}^k d_{p_j}\, \Big( {x(p_j)^m\,W^{(g)}_{k}(p_{K})\over dx(p_j)}  \Big)\cr
\eea
due to theorem \ref{thPkpol}.
}

\vs

{\bf Theorem \ref{thintPhi}} \proof{
The case $k=1$, $g=0$ is easy:
\beq
\Res_{p_2\to p_1} \Phi(p) B(p_1,p_2) = d\Phi(p_2) = y(p_1)dx(p_1).
\eeq
We prove the theorem by recursion on $g$ and $k$. Suppose it is proved for all $k'$ for $g'\leq g-1$, and for $k'\leq k-1$ if $g'=g$.
We write $K=\{1,\dots,k\}$ and $K'=\{1,\dots,k-1\}$.
Then we have from \eq{defWkgrecursive}:
\bea
&& \Res_{p_{k}\to \bfa} \Phi(p_{k}) W_{k+1}^{(g)}(p,p_1,\dots,p_k) \cr
&=& \Res_{p_{k}\to \bfa} \Res_{q\to \bfa} \Phi(p_{k}) \,{dE_{q}(p)\over \om(q)}\,\Big(
\qquad \sum_{m=0}^g \sum_{J\subset K'} W_{j+2}^{(m)}(q,p_J,p_k)W_{k-j}^{(g-m)}(\qbar,p_{K'/J})  \cr
&& \qquad +\sum_{m=0}^g \sum_{J\subset K'} W_{j+1}^{(m)}(q,p_J)W_{k-j+1}^{(g-m)}(\qbar,p_{K'/J},p_k)
+ W_{k+2}^{(g-1)}(q,\qbar,p_{K'},p_k) \Big) .\cr
\eea
Then we exchange the contours of integration:
\beq
\Res_{p_{k}\to \bfa} \Res_{q\to \bfa} = \Res_{q\to \bfa} \Res_{p_{k}\to \bfa} + \Res_{q\to \bfa} \Res_{p_{k}\to q,\qbar}.
\eeq
Thus
\bea
&& \Res_{p_{k}\to \bfa} \Phi(p_{k}) W_{k+1}^{(g)}(p,p_1,\dots,p_k) \cr
&=& \Res_{q\to \bfa} \Res_{p_{k}\to \bfa} \Phi(p_{k}) \,{dE_{q}(p)\over \om(q)}\,\Big(
\qquad \sum_{m=0}^g \sum_{J\subset K'} W_{j+2}^{(m)}(q,p_J,p_k)W_{k-j}^{(g-m)}(\qbar,p_{K'/J})  \cr
&& \qquad +\sum_{m=0}^g \sum_{J\subset K'} W_{j+1}^{(m)}(q,p_J)W_{k-j+1}^{(g-m)}(\qbar,p_{K'/J},p_k)
+ W_{k+2}^{(g-1)}(q,\qbar,p_{K'},p_k) \Big) \cr
&& + \Res_{q\to \bfa} \Res_{p_{k}\to q,\qbar} \Phi(p_{k}) \,{dE_{q}(p)\over \om(q)}\,\Big(
\qquad \sum_{m=0}^g \sum_{J\subset K'} W_{j+2}^{(m)}(q,p_J,p_k)W_{k-j}^{(g-m)}(\qbar,p_{K'/J})  \cr
&& \qquad +\sum_{m=0}^g \sum_{J\subset K'} W_{j+1}^{(m)}(q,p_J)W_{k-j+1}^{(g-m)}(\qbar,p_{K'/J},p_k)
+ W_{k+2}^{(g-1)}(q,\qbar,p_{K'},p_k) \Big) .\cr
\eea
The first term is computed from the recursion hypothesis, and the second term can exist only if the correlation function containing $p_k$ has poles at $p_k=q$ or $p_k=\qbar$, and  from theorem \ref{thpolesWkgbp},
this can happen only if the correlation function containing $p_k$  is a Bergmann kernel. That gives:
\beq
\begin{array}{l}
 \Res_{p_{k}\to \bfa} \Phi(p_{k}) W_{k+1}^{(g)}(p,p_1,\dots,p_k) \cr
=  \Res_{q\to \bfa} {dE_{q}(p)\over \om(q)}\,\Big(
\qquad \sum_{m=0}^g \sum_{J\subset K'} (2m+(j+1)-2) W_{j+1}^{(m)}(q,p_J) W_{k-j}^{(g-m)}(\qbar,p_{K'/J})  \cr
 \qquad +\sum_{m=0}^g \sum_{J\subset K'} (2(g-m)+(k-j)-2) W_{j+1}^{(m)}(q,p_J) W_{k-j}^{(g-m)}(\qbar,p_{K'/J}) \cr
 \qquad + (2(g-1)+k+1-2) W_{k+1}^{(g-1)}(q,\qbar,p_{K'}) \Big) \cr
 + \Res_{q\to \bfa} \Res_{p_{k}\to q,\qbar} \Phi(p_{k}) \,{dE_{q}(p)\over \om(q)}\,\Big(
\qquad  B(q,p_k)W_{k+1}^{(g)}(\qbar,p_{K'})  +  W_{k}^{(g)}(q,p_{K'}) B(\qbar,p_k) \Big) \cr
=  (2g+k-3)\, \Res_{q\to \bfa} {dE_{q}(p)\over \om(q)}\,\Big( \sum_{m=0}^g \sum_{J\subset K'}  W_{j+1}^{(m)}(q,p_J) W_{k-j}^{(g-m)}(\qbar,p_{K'/J})  \cr
\qquad +  W_{k+1}^{(g-1)}(q,\qbar,p_{K'}) \Big) \cr
 + \Res_{q\to \bfa} {dE_{q}(p)\over y(q)-y(\qbar)}\,\Big( y(q) W_{k+1}^{(g)}(\qbar,p_{K'})  +  y(\qbar) W_{k}^{(g)}(q,p_{K'})  \Big) \cr
=  (2g+k-3)\, W_k(p,p_1\dots,p_{k-1}) \cr
 + \Res_{q\to \bfa} {dE_{q}(p)\over y(q)-y(\qbar)}\,\Big( y(q) (W_{k+1}^{(g)}(\qbar,p_{K'})+W_{k}^{(g)}(q,p_{K'}))  +  (y(\qbar)-y(q)) W_{k}^{(g)}(q,p_{K'})  \Big) \cr
=  (2g+k-3)\, W_k(p,p_1\dots,p_{k-1}) \cr
 + \Res_{q\to \bfa} {dE_{q}(p)\over y(q)-y(\qbar)}\, (y(\qbar)-y(q)) W_{k}^{(g)}(q,p_{K'})   \cr
=  (2g+k-3)\, W_k(p,p_1\dots,p_{k-1})  - \Res_{q\to \bfa} dE_{q}(p)\, W_{k}^{(g)}(q,p_{K'})   \cr
=  (2g+k-3)\, W_k(p,p_1\dots,p_{k-1})  + {1\over 2}\,\Res_{q\to \bfa} dS_{q,o}(p) \, W_{k}^{(g)}(q,p_{K'}) \cr
 \qquad  -  {1\over 2}\,\Res_{q\to \bfa} dS_{\qbar,o}(p) \, W_{k}^{(g)}(q,p_{K'})  \cr
=  (2g+k-3)\, W_k(p,p_1\dots,p_{k-1})  + {1\over 2}\,\Res_{q\to \bfa} dS_{q,o}(p) \, W_{k}^{(g)}(q,p_{K'}) \cr
 \qquad -  {1\over 2}\,\Res_{q\to \bfa} dS_{q,o}(p) \, W_{k}^{(g)}(\qbar,p_{K'})  \cr
=  (2g+k-3)\, W_k(p,p_1\dots,p_{k-1})  + \Res_{q\to \bfa} dS_{q,o}(p) \, W_{k}^{(g)}(q,p_{K'}) \cr
=  (2g+k-3)\, W_k(p,p_1\dots,p_{k-1})  - \Res_{q\to p} dS_{q,o}(p) \, W_{k}^{(g)}(q,p_{K'}) \cr
 -{1\over 2i\pi} \sum_j \left( \oint_{\underline{\bcycle}_j} B(q,p) \oint_{\underline{\acycle}_j} W_{k}^{(g)}(q,p_{K'}) - \oint_{\underline{\acycle}_j} B(q,p) \oint_{\underline{\bcycle}_j} W_{k}^{(g)}(q,p_{K'})\right) \cr
=  (2g+k-3)\, W_k(p,p_1\dots,p_{k-1})  - \Res_{q\to p} dS_{q,o}(p) \, W_{k}^{(g)}(q,p_{K'}) \cr
 - du^t(p)\,\left( (1+\kappa\tau) \oint_{\underline{\acycle}} W_{k}^{(g)}(q,p_{K'}) - \kappa \oint_{\underline{\bcycle}} W_{k}^{(g)}(q,p_{K'})\right) \cr
=  (2g+k-3)\, W_k(p,p_1\dots,p_{k-1})  - \Res_{q\to p} dS_{q,o}(p) \, W_{k}^{(g)}(q,p_{K'}) \cr
 \qquad - du^t(p)\,\oint_{\acycle} W_{k}^{(g)}(q,p_{K'})  \cr
=  (2g+k-3)\, W_k(p,p_1\dots,p_{k-1})  - \Res_{q\to p} dS_{q,o}(p) \, W_{k}^{(g)}(q,p_{K'})   \cr
=  (2g+k-2)\, W_k(p,p_1\dots,p_{k-1})
\end{array}
\eeq

}



